\chapter{糖类\texorpdfstring{\quad}{ }蛋白质}
在自然界里，糖类、蛋白质和\cref{chp:hydrocarbon_derivatives}学过的脂肪都是动植物等进行生命活动的重要有机物。
糖类是绿色植物通过光合作用吸收日光能的产物，它把能量储存起来，又作为动、植物所需要能量的重要来源。
蛋白质是组成各种细胞的基础物质。
在这章里，我们将学习这两类物质。

\section{单糖}
葡萄糖、蔗糖、淀粉和纤维素等都属于糖类。
糖类也叫做碳水化合物。
碳水化合物这个名称的得来，是由于最先发现的这一类化合物都是由 \ce{C}、\ce{H}、\ce{O} 三种元素组成的，分子中氢原子和氧原子数之比恰好是 $2:1$。
糖类的化学式可以用通式 \ce{C_$n$(H2O)_$m$} 来表示（$n$ 和 $m$ 可以相同，也可以不同）。
因此就错误地把它们看做是碳的水化物了。
事实上，碳水化合物这个名称并不能反映它们的结构特点。
首先在碳水化合物分子中，\ce{H} 和 \ce{O} 并不是以结合成水的形式存在着；还有，已经发现许多碳水化合物，在它们的分子中，\ce{H} 和 \ce{O} 的比并不等于 $2:1$，如鼠李糖 \ce{C6H12O5}；而且，有许多符合 \ce{C_$n$(H2O)_$m$} 通式的物质也并不属于碳水化合物，如甲醛 \ce{CH2O}、乙酸 \ce{C2H4O2} 等。
因此，碳水化合物这个名称虽然仍旧沿用，但早已失去原来的意义。
\emph{从结构上看，糖类一般是多羟基醛或多羟基酮，以及能水解生成它们的物质}。

根据分子的结构，糖类可以分为单糖、低聚糖和多糖。
在这节里先学习单糖。

单糖中最重要的是葡萄糖和果糖。
\subsection{葡萄糖}
葡萄糖的化学式是 \ce{C6H12O6}。
葡萄糖是自然界中分布最广的单糖。
葡萄糖存在于葡萄汁和其它带甜味的水果里。
蜂蜜里也含有葡萄糖。
葡萄糖不如蔗糖甜。

\subsubsection{葡萄糖的性质和结构}
葡萄糖是白色晶体，溶于水。
葡萄糖在水溶液里的一些性质表明它的结构简式是：
\[ \ce{CH2OH-CHOH-CHOH-CHOH-CHOH-CHO } \]

葡萄糖是一种多羟基醛。
\begin{Experiment}
在一支洁净的试管里加入 2\% 硝酸银溶液 \qty{2}{ml}，振荡试管，同时逐滴加入 2\% 稀氨水，直到析出的沉淀恰好溶解为止。所得澄清溶液就是银氨溶液。在盛着银氨溶液的试管里加入 10\% 葡萄糖溶液 \qty{1}{ml}，把试管放在水浴里加热 \qty{3}{5}{min}，观察银镜的出现。

在另一支试管里加入 10\% 氢氧化钠溶液 \qty{2}{ml}，滴加入 5\% 硫酸铜溶液 \numrange{4}{5} 滴。可以观察到淡蓝色氢氧化铜沉淀生成。立即加入 10\% 葡萄糖溶液 \qty{2}{ml}。加热后，可以观察到有红色氧化亚铜沉淀生成。
\end{Experiment}
根据上述实验，葡萄糖具有的醛基，被氧化成为羧基，葡萄糖变为葡萄糖酸（简称葡糖酸）。这两个反应可以用化学方程式简单表示如下：
\begin{multline*}
  \ce{CH2OH-(CHOH)4-CHO + 2[Ag(NH3)2]+ + 2OH- -> } \\ 
  \schemestart
  \chemname{\ce{CH2OH-(CHOH)4-COOH}}{\footnotesize 葡萄糖酸}
  \+ \ce{2Ag v + H2O + 4NH3}
  \schemestop\chemnameinit{}
\end{multline*}
\begin{multline*}
  \ce{CH2OH-(CHOH)4-CHO + 2Cu(OH)2 -> } \\ \ce{ CH2OH-(CHOH)4-COOH + CuO v + 2H2O }
\end{multline*}

葡萄糖具有醛基，象醛类一样，能被还原剂还原。
葡萄糖加氢时还原为六元醇。
葡萄糖具有醇羟基，能跟酸起酯化反应。

\subsubsection{葡萄糖的制法}
在工业上，葡萄糖通常是采用淀粉作原料，用硫酸等无机酸作催化剂，起水解反应而制得的。
% \[ \ce{(C6H10O5)_$n$ + $n$H2O ->[\text{催化剂}] $n$C6H12O6 } \]
\[\schemestart
\chemname{\ce{(C6H10O5)_$n$}}{\footnotesize 淀粉} \+ $n$\,\ce{H2O}
\arrow(.mid east--.mid west){->[\footnotesize 催化剂]}
$n$\,\chemname{\ce{C6H12O6}}{\footnotesize 葡萄糖}
\schemestop\chemnameinit{}\]

\subsubsection{葡萄糖的用途}
葡萄糖是一种重要的营养物质，它是人类生命活动所需能量的重要来源之一，它在人体组织中进行氧化反应，放出热量，以供应人们所需要的能。
\[ \ce{ C6H12O6\,\text{（固）} + 6O2\,\text{（气）} -> 6CO2\,\text{（气）} + 6H2O\,\text{（液）}} + \qty{669.9}{kCal} \]

葡萄糖用于医药，也用于制造糖果。
制镜工业和热水瓶胆镀银常用葡萄糖作还原剂。

\subsection{果糖}
果糖的化学式也是 \ce{C6H12O6}。
果糖在自然界中分布也较广，存在于水果中，也存在于蜂蜜中。
果糖比蔗糖甜。
纯净的果糖是白色晶体，它不易结晶，通常是粘稠的液体，易溶于水。
实验证明，果糖的结构简式是:
\[ \ce{ CH2OH-CHOH-CHOH-CHOH-CO-CH2OH }\]

因此，果糖是一种多羟基酮。

\begin{Reading}[补充阅读：]{核糖}
核糖是个含有五个碳原子的糖——戊糖。
核糖是细胞核的重要组成部分，是人类生命活动中不可缺少的物质。
重要的核糖有核糖和脱氧核糖。
它们的结构简式是：
\begin{align*}
      \text{核糖}&\quad \ce{ CH2OH-CHOH-CHOH-CHOH-CHO }\\
  \text{脱氧核糖}&\quad \ce{ CH2OH-CHOH-CHOH-CH2-CHO } \\
\end{align*}
\end{Reading}

\begin{Practice}[习题]
\begin{question}
  \item 葡萄糖的式量是 180，其中含碳 40\%，氢 6.7\%，其余是氧。求葡萄糖的化学式。
  \item 写出下列化学方程式：
  \begin{tasks}[label=(\arabic*)]
    \task 葡萄糖被银氨溶液氧化，
    \task 葡萄糖加氢还原为六元醇。
  \end{tasks}
  \item 葡萄糖跟果糖的分子结构有什么区别？
  \item 患有糖尿病的人，尿中含有糖分（葡萄糖）较多，怎样检验一个病人是否患有糖尿病？
  \item 热水瓶胆要镀银通常是用葡萄糖为还原剂，如每一个热水瓶胆要镀银 \qty{0.3}{g}，一天生产 5000 个这种热水瓶胆的工厂，要用多少千克 98\% 的葡萄糖来还原银氨溶液。
\end{question}
\end{Practice}

\section{二糖}
糖类水解后能生成几个分子单糖的叫做\Concept{低聚糖}。
根据水解后生成的单糖分子是二个、三个等，低聚糖又分为二糖、三糖等。
其中最重要的是二糖。
蔗糖和麦芽糖都是二糖。

\subsection{蔗糖}
二糖里人们最熟悉的是蔗糖，化学式是 \ce{C12H22O11}。
蔗糖是无色晶体，溶于水。
蔗糖是重要的甜味食物，存在于不少植物体内，以甘蔗（含糖 11\%～17\%）和甜菜（含糖 14\%～26\%）的含量最多。
\begin{Experiment}
在两个洁净的试管里，各加入 10\% 的蔗糖溶液 \qty{1}{ml}，在一个试管加入 3 滴稀硫酸（$1:5$），把两个试管都放在水浴上加热，然后在两个试管里各加入 \qty{2}{ml} 银氨溶液。对比两个试管里的反应，观察有没有银镜反应发生。

用硫酸铜和氢氧化钠溶液代替银氨溶液作上述实验。
\end{Experiment}
从上述实验可以看出，蔗糖不发生银镜反应，这因为它的分子结构中不含有醛基。
蔗糖不显还原性，是一种非还原糖。
在硫酸的催化作用下，蔗糖水解生成一分子葡萄糖和一分子果糖：
\[\schemestart
\ce{ C12H22O11 + H2O }
\arrow(.mid east--.mid west){->[\scriptsize 催化剂]}
\chemname{\ce{C6H12O6}}{\footnotesize 葡萄糖} \+
\chemname{\ce{C6H12O6}}{\footnotesize 果糖} 
\schemestop\chemnameinit{}\]
因此蔗糖水解后能发生银镜反应，也能还原氢氧化铜。

\subsection{麦芽糖}
麦芽糖是白色晶体（常见的麦芽糖是没有结晶的糖膏），易溶于水，有甜味，但不如蔗糖甜。
麦芽糖能发生银镜反应，因为它的分子结构中还含有醛基。
麦芽糖是一种还原糖。
在硫酸等的催化作用下，麦芽糖发生水解反应，生成两分子葡萄糖。
\[\schemestart
\chemname{\ce{C12H22O11}}{\footnotesize 麦芽糖} \+ \ce{2H2O}
\arrow(.mid east--.mid west){->[\scriptsize 催化剂]}
\chemname{\ce{2C6H12O6}}{\footnotesize 葡萄糖} 
\schemestop\chemnameinit{}\]

从麦芽糖和蔗糖的水解反应可以看出，二糖水解后能够生成两分子的单糖，而\Concept{单糖}是不能水解的简单糖类。

麦芽糖是用含淀粉较多的农产品如大米、玉米、薯类等作为原料，在淀粉酶（大麦芽产生的酶）的作用下，在约 \qty{60}{\celsius} 时，发生水解反应而生成的：
\[\schemestart
\chemname{\ce{2(C6H10O5)_$n$}}{\footnotesize 淀粉} \+ \ce{$n$H2O}
\arrow(.mid east--.mid west){->[\scriptsize 催化剂]}
$n$\,\chemname{\ce{C12H22O11}}{\footnotesize 麦芽糖}
\schemestop\chemnameinit{}\]
麦芽糖也用作甜味食物。

\begin{Practice}[习题]
\begin{question}
  \item 怎样鉴别麦芽糖和蔗糖？
  \item 怎样使二糖转化为单糖？怎样用实验证明蔗糖已转化为葡萄糖？
  \item 在蔗糖里加入浓硫酸，会发生什么现象？为什么？
  \item \qty{2}{mol} 蔗糖水解生成的果糖和葡萄糖各几克？
\end{question}
\end{Practice}

\section{多糖}
\Concept{多糖}是由很多个单糖分子按照一定的方式，通过在分子间脱去水分子，结合而成的。
多糖在性质上跟单糖、低聚糖不同，一般不溶于水，没有甜味，没有还原性。
淀粉和纤维素是最重要的多糖,它们的通式都是 \ce{(C6H10O5)_$n$}。
淀粉和纤维素的分子里所包含的单糖单元 \ce{(C6H10O5)} 的数目不同，即 $n$ 值不同。
淀粉和纤维素在结构上也有所不同。

\subsection{淀粉}
淀粉是绿色植物进行光合作用的产物，主要存在于植物的种子或块根里。
谷类中含淀粉较多。
例如，大米约含淀粉 80\%，小麦约含 70\%，马铃薯约含 20\% 等。

淀粉本身主要是由两部分组成——直链淀粉（约占淀粉的 20\%）和支链淀粉（约占淀粉的 80\%）组成。

直链淀粉分子中约含有几百个葡萄糖单元，它的式量从几万到十几万；支链淀粉分子中约含有几千个葡萄糖单元，它的式量约为几十万。
淀粉是一类式量很大的化合物。
\emph{这类式量很大的化合物叫}\Concept{高分子化合物}。

\subsubsection{淀粉的性质和结构}
淀粉是一种白色粉末状的物质。
直链淀粉能溶于热水。 
支链淀粉不溶于水，但能在水中胀大而湿润。
\begin{Experiment}
在试管 1 和试管 2 里各放入 \qty{0.5}{g} 淀粉，在试管 1 里加入 \qty{4}{ml} 20\% 硫酸溶液，在试管 2 里加入 \qty{4}{ml}水，都加热 \qty{3}{4}{min}。用碱液中和试管 1 里的硫酸溶液，把一部分液体倒入试管 3。在试管 2 和试管 3 里都加入碘溶液，观察有没有蓝色出现。在试管 1 里加入银氨溶液，稍加热后，观察试管内壁上有无银镜出现。
\end{Experiment}
从上述实验，知道淀粉经用酸作催化剂发生水解后，生成还原性的单糖，能起银镜反应。

不论是直链淀粉，还是支链淀粉，在稀酸作用下都发生水解，生成一系列的产物，最后得到葡萄糖。
% \[ \ce{ (C6H10O5)_$n$  +  $n$H2O ->[\text{催化剂}] $n$C6H12O6 } \]
\[\schemestart
\chemname{\ce{(C6H10O5)_$n$}}{\footnotesize 淀粉}
\+ \ce{$n$\,H2O}
\arrow(.mid east--.mid west){->[\scriptsize 催化剂]}
$n$\,\chemname{\ce{C6H12O6}}{\footnotesize 葡萄糖}
\schemestop\chemnameinit{}\]

淀粉在人体内也进行水解。
人们在咀嚼食物的时候，淀粉受唾液所含淀粉酶（一种蛋白质，能对淀粉的水解起催化作用）的作用，开始水解。
我们在吃饭时多加咀嚼，就会感到甜味。

淀粉在小肠里，在胰脏分泌出的淀粉酶的作用下，继续进行水解。
生成的葡萄糖经过肠壁的吸收，进入血液，供人体组织的营养需要。

\subsubsection{淀粉的用途}
淀粉是食物的一种重要成分。
它也是一种工业原料，可以用来制葡萄糖和酒精等。
淀粉在淀粉酶的作用下，先转化为麦芽糖，再转化为葡萄糖。
葡萄糖受到酒曲里的酒化酶的作用，变为酒精。
这就是含淀粉物质酿酒的主要过程。
葡萄糖变为酒精的反应可以简略表示如下：
\[\ce{ C6H12O6 ->[\text{催化剂}] 2C2H5OH + 2CO2 }\]

\subsection{纤维素}
纤维素是构成细胞壁的基础物质。
木材约有一半是纤维素。
棉花是自然界中较纯粹的纤维素，约含纤维素 92\%～95\%。
脱脂棉和无灰滤纸差不多是纯粹的纤维素。

纤维素分子中大约含有几千个葡萄糖单元，它的式量约为几十万。

\subsubsection{纤维素的性质和结构}
纤维素是白色、无臭、无味的物质，不溶于水，也不溶于一般的有机溶剂。

跟淀粉一样，纤维素不显示还原性。
纤维素可以发生水解，但较淀粉困难。
纤维素在稀酸和一定压强下长时间加热，可发生水解，水解的最后产物是葡萄糖。
\begin{Experiment}
把少许棉花或几片碎滤纸放入试管里，加入 70\% 硫酸 \qtyrange{3}{4}{ml}，用玻璃棒把棉花捣烂，形成无色粘稠液体。
把这试管放在水浴中加热，约 \qty{15}{min}，放冷后倾入盛有 \qty{20}{ml} 水的烧杯里，用氢氧化钠溶液中和。
取出一部分已中和的溶液，用硫酸铜和氢氧化钠溶液作试剂试验，可以观察到有红色氧化亚铜生成。
\end{Experiment}
纤维素的水解反应可以表示如下：
% \[ \ce{ (C6H10O5)_$n$ + $n$H2O ->[\text{催化剂}] $n$C6H12O6 }\]
\[\schemestart 
\chemname{\ce{(C6H10O5)_$n$}}{\footnotesize 纤维素} \+ \ce{$n$\,H2O}
\arrow(.mid east--.mid west){->[\scriptsize 催化剂]}
\chemname{\ce{$n$C6H12O6}}{\footnotesize 葡萄糖}
\schemestop\chemnameinit{}\]

纤维素分子是由很多个葡萄糖单元构成的。
每一个葡萄糖单元有三个醇羟基，因此，纤维素分子也可用 \ce{[C6H7O2(OH)3]_$n$} 表示。
由于醇羟基的存在，所以纤维素能够表现出醇的一些性质，如生成硝酸酯、乙酸酯等。

\subsubsection{纤维素的用途}
纤维素常用于制造纤维素硝酸酯、纤维素乙酸酯、粘胶纤维和造纸等等。

\paragraph{制造纤维素硝酸酯}\mbox{}\par
纤维素硝酸酯俗名硝酸纤维，是由棉花(成分是纤维素) 跟浓硝酸、浓硫酸的混和物在一定条件下反应制得的。
\begin{Experiment}
把 \qty{5}{ml} 浓硝酸（密度 \qty{1.4}{g/cm^3}）和 \qty{10}{ml} 浓硫酸（密度 \qty{1.84}{g/cm^3}）先后加入烧杯里配好混和液。放入一小块棉花，浸 \qtyrange{8}{10}{mini}，然后取出洗净，晾干，把晾干的纤维素硝酸酯和一小块棉花同时点燃，比较它们的燃烧现象。
\end{Experiment}
可以看出纤维素硝酸酯的燃烧比棉花迅速得多。生成纤维素硝酸酯的反应表示如下：
\begin{multline*}
  \schemestart 
  \chemname{$\left[\chemfig{{(C_6H_7O_2)}-[,,,,draw=none]{\phantom{C}}(-[1,1.4]OH)(-[7,1.4]OH)-OH}\right]_n$}{\footnotesize 纤维素}
\+ 
$3n$\,\chemname{\ce{HO-NO2}}{\footnotesize 硝酸}
\arrow(.mid east--.mid west){->[\scriptsize 浓硫酸]} \schemestop\\
\schemestart 
\chemname{$\left[\chemfig{{(C_6H_7O_2)}-[,,,,draw=none]{\phantom{C}}(-[1,1.4]O-[:0]NO_2)(-[7,1.4]O-[:0]NO_2)-O-[:0]NO_2}\right]_n$}{\footnotesize 纤维素三硝酸酯}
\+
\ce{3$n$H2O}
\schemestop\chemnameinit{}\end{multline*}

纤维素一般不容易完全酯化而生成三硝酸酯（含 \ce{N} 14.14\%）。
酯化程度不同的硝酸酯在性质上也有所不同，工业上把含氮量低的纤维素硝酸酯叫做胶棉（含 \ce{N} 10.5\%～12\%），含氮量高的叫做火棉（含 \ce{N} 12.5\%～13.8\%）。

火棉在外表上跟棉花相似，遇火即迅速燃烧，在密闭容器中发生爆炸，可用作无烟火药。

胶棉也易于燃烧，但并不爆炸。
胶棉用于制喷漆。
胶棉的乙醇—乙醚溶液俗名珂璆酊，用于封瓶口等。

\paragraph{制造纤维素乙酸酯}\mbox{}\par
纤维素乙酸酯俗名醋酸纤维，是由棉花跟乙酸—乙酸酐（\ce{(CH3CO)2O}）的混和物在一定条件下反应制得的。

纤维素乙酸酯多用于制电影胶片片基。它的优点是不易着火。

\paragraph{制造粘胶纤维}\mbox{}\par
把纤维素依次用氢氧化钠浓溶液和二硫化碳处理，再把生成物溶解于氢氧化钠稀溶液中即形成粘胶液。
把粘胶液经过细孔压入稀酸溶液中，重新生成纤维素，即是粘胶纤维。
粘胶纤维中的长纤维俗称人造丝，短纤维俗称人造棉，都可供纺织用。
如果把粘胶液通过狭缝压入稀酸中，可制成透明的薄膜，俗称玻璃纸。

\paragraph{造纸}\mbox{}\par
造纸是我国古代重大发明之一，是我国人民对人类文化的伟大贡献。
一九五七年，在陕西省西安市灞桥西汉前期的汉墓中发现了西汉的早期纸张（由麻类植物纤维所制）。
这说明早在公元前二世纪在我国已使用由植物纤维制造的纸张了。
近代造纸的原料主要是用植物纤维一木材纤维和草本植物纤维，象芦苇、稻草、麦秸、蔗渣等。
为了把植物纤维制成纸浆，可以用机械法（磨碎）和化学处理法。
化学处理法是用亚硫酸氢钙或氢氧化钠等化学药品使纤维素原料中的非纤维素成分溶解除去，使纤维素分离出来，得到化学纸浆。
制成纸浆后，再经过漂白、打浆、抄纸（铺成薄层）、烘干而成纸。

\begin{Practice}[习题]
\begin{question}
  \item 在以淀粉为原料生产葡萄糖的水解过程中，用什么方法来检验淀粉已完全水解？
  \item 某厂用含淀粉 54\% 的薯干 \qty{2}{t} 来制取酒精。如果在发酵过程中有 85\% 的淀粉转化为酒精，制得的酒精又含水 5\%，可得这样的酒精多少吨？
  \item 没有成熟的苹果肉遇碘显蓝色，成熟的苹果汁能还原银氨溶液。怎样解释这两种现象？
  \item 要制取 \qty{1}{t} 纤维素三硝酸酯，需精制过的棉短绒和 60\% 的浓硝酸各多少？
  \item 试举例说明单糖、二糖和多糖之间相互转化的关系。
\end{question}
\end{Practice}

\section{氨基酸}
\subsection{氨基酸}
羧酸分子里烃基上的氢原子被氨基取代后的生成物，叫做氨基酸。下面是几种氨基酸的名称和结构式：

甘氨酸（氨基乙酸）
\[\chemfig{CH_2(-[:-90]NH_2)-COOH}\]

丙氨酸（$\alpha$-氨基丙酸）
\[\chemfig{CH_3-CH(-[:-90]NH_2)-COOH}\]

苯丙氨酸（$\alpha$-氨基-$\beta$-苯基丙酸）
\[\chemfig{[:-30]*6(-=-(-CH_2-CH(-[:-90]NH_2)-COOH)=-=)}\]

谷氨酸（$\alpha$-氨基戊二酸）
\[\chemfig{HOOC-CH_2-CH_2-CH(-[:-90]NH_2)-COOH}\]

上述氨基酸都是 $\alpha$-氨基酸，即羧酸分子里的 $\alpha$ 氢原子（即离羧基最近的碳原子上的氢原子，离羧基次近的碳原子上的氢原子依次为 $\beta$ 氢原子等）被氨基取代的生成物。

氨基酸是晶体，熔点较高，一般在 \qtyrange{200}{300}{\celsius}。氨基酸一般可溶于水，但不溶于乙醚。

氨基酸分子里含有羧基和氨基，它们既是酸、又是碱，因此，氨基酸跟酸或碱作用都可以生成盐。例如：
\[\schemestart
\chemfig{CH_2(-[:-90]NH_2)-COOH} \+ \ce{HCl}  \arrow(.mid east--.mid west)
\chemfig{\vphantom{CH_2}-[@{op,.7},,,,draw=none]CH_2(-[:-90]NH_3^+)-COOH-[@{cl,.5},,,,draw=none]Cl^{-}}
\polymerdelim[delimiters ={[]},height=3pt,depth =30pt,indice={}]{op}{cl}
\schemestop\]
\[\schemestart
\chemfig{CH_2(-[:-90]NH_2)-COOH} \+ \ce{NaOH}  \arrow(.mid east--.mid west)
\chemfig{\vphantom{CH_2}-[@{op,.7},,,,draw=none]CH_2(-[:-90]NH_2)-COO^{-}-[@{cl,.5},,,,draw=none]Na^{+}}
\polymerdelim[delimiters ={[]},height=3pt,depth =30pt,indice={}]{op}{cl}
\+ \ce{H2O}
\schemestop\]
前一种盐具有酸性（因为分子中还含有羧基），后一种盐具有碱性（因为分子中还含有氨基）。

蛋白质在酸、碱或酶的作用下发生水解，经过多肽，最后得到多种 $\alpha$-氨基酸。
因此，氨基酸是构成蛋白质的“基石”。

\subsection{多肽}
一分子氨基酸中的羧基跟另一分子氨基酸中的氨基之间消去水分子，经缩合反应而生成的产物叫做肽，其中的酰胺基结构 \ce{-CO-NH-} 叫做\Concept{肽键}。
两个氨基酸分子起反应生成肽的化学方程式如下：
\begin{multline*} 
  \schemestart
  \chemfig{H_2N-CH(-[6]R)-CO-OH}\tikz[remember picture]\coordinate(a); \+ \tikz[remember picture]\coordinate(b);\chemfig{H-NH-CH(-[6]R)-COOH}
  \arrow(.mid east--.mid west)
  \schemestop\\
  \schemestart
  \tikz[remember picture]\coordinate(c);\chemfig{H_2N-CH(-[6]R)-CO-NH-CH(-[6]R)-COOH}
  \+ \ce{H2O}
  \schemestop
  \tikz[remember picture,overlay]{
    \draw[densely dashed]([xshift=-2.1em,yshift=-1ex]a)rectangle([xshift=1.4em,yshift=3ex]b);
    \draw[decorate,decoration={brace,mirror,amplitude=5pt}]([xshift=5.5em,yshift=-0.5ex]c)--++(5em,0)node[midway,below=5pt]{\footnotesize 肽键};
  }
\end{multline*}

由两个氨基酸分子消去水分子而形成含有一个肽键的化合物是二肽。
由多个氨基酸分子消去水分子而形成含有多个肽键的化合物是\Concept{多肽}。
多肽通常呈链状。
蛋白质水解得到多肽；多肽进一步水解，最后得到 $\alpha$-氨基酸。
多肽和蛋白质之间没有严格的区别，一般常把式量小于 \num{10000} 的叫做多肽。

\section{蛋白质}
\subsection{蛋白质}
蛋白质广泛存在于生物体内，是组成细胞的基础物质。
动物的肌肉、皮肤、血液、乳汁以及发、毛、角、蹄的主要成分都是蛋白质。
植物的各种器官也都含有蛋白质，例如，小麦的种子里约含 18\% 蛋白质。
所有的酶都是蛋白质。
滤过性病毒含有蛋白质，但它的成分更复杂。

蛋白质是由 $\alpha$-氨基酸通过肽键构成的高分子化合物。
式量相差很大，有的几万、几十万，个别的甚至上千万。
核蛋白的式量更大，如烟草斑纹病毒的核蛋白的式量超过二千万，腺病毒的核蛋白的式量还要大。

\begin{Reading}[补充阅读]{}
核蛋白通常由蛋白质和核酸组成。
核酸有两种，含核糖的是核糖核酸（RNA），含脱氧核糖的是脱氧核糖核酸 （DNA）。脱氧核糖核酸呈双螺旋结构（参见彩图），以氢键相结合，它是跟生物遗传密切有关的物质。
\end{Reading}

淀粉、纤维素、蛋白质等都是自然界中存在的天然有机高分子化合物。

蛋白质的结构非常复杂。
多肽链内多种 $\alpha$-氨基酸以一定顺序排列，多肽链跟多肽链之间以氢键相结合，呈空间结构。
不同结构的蛋白质的生理功能不同，即便是化学组成不变，而只是空间结构发生了改变，它的生理功能也会发生变化。

蛋白质除含碳、氢、氧外，都含氮和少量硫，含氮量为 15\%～18\%，含硫量为 0.3\%～2.5\%。

下面简单介绍蛋白质的一些性质：
\subsubsection{两性}
蛋白质是由 $\alpha$-氨基酸通过肽键构成的高分子化合物，在蛋白质分子中存在着氨基和羧基。
因此跟氨基酸相似，蛋白质也是两性物质。
\subsubsection{盐析}
\begin{Experiment}
在盛有鸡蛋白溶液的试管里，缓慢地加入饱和硫酸铵或硫酸钠溶液，可以观察到有沉淀析出。把少量带有沉淀的液体加入盛有清水的试管里，观察沉淀是否溶解。
\end{Experiment}
有些蛋白质如鸡蛋白能够溶解在水里形成溶液。
蛋白质的分子的直径达到了胶体微粒的大小，所以，蛋白质溶液具有胶体的性质。

少量的盐（如硫酸铵、硫酸钠等）能促进蛋白质的溶解，但如向蛋白质水溶液中，加入浓的无机盐溶液，可使蛋白质的溶解度降低，而从溶液中析出。这种作用叫做盐析。
这样析出的蛋白质仍旧可以溶解在水中，而不影响原来蛋白质的性质。
因此，盐析是个可逆过程。
利用这个性质，可以采用多次盐析的方法来分离、提纯蛋白质。

\subsubsection{变性}
\begin{Experiment}
在两个试管里各盛 \qty{3}{ml} 鸡蛋白溶液，给一个试管加热，向另一个试管加入少量乙酸铅溶液，观察发生的现象。把凝结的蛋白和生成的沉淀分别放入两个盛清水的试管里，观察是否溶解。 
\end{Experiment}

在热、酸、碱、重金属盐、紫外线等作用下，蛋白质会发生性质上的改变而凝结起来。
这种凝结是不可逆的，不能再使它们恢复成为原来的蛋白质。
蛋白质的这种变化叫做变性。
蛋白质变性后，就丧失了原有的可溶性，并且失去了它们生理上的作用。
高温消毒灭菌就是利用加热使蛋白质凝固从而使细胞死亡。
重金属盐（如铜盐、铅盐、汞盐等）能使蛋白质凝结，所以会使人中毒。

\subsubsection{颜色反应}
\begin{Experiment}
在盛有 \qty{2}{ml} 鸡蛋白溶液的试管里，滴入几滴浓硝酸，加微热，观察析出的沉淀的颜色。
\end{Experiment}

蛋白质可以跟许多试剂发生颜色反应。
例如，有些蛋白质跟浓硝酸作用时呈黄色。
有这种反应的蛋白质分子中一般有苯环存在，如肽链中含有苯丙氨酸单元等。
在使用浓硝酸时，不慎溅在皮肤上而使皮肤呈现黄色，就是由于浓硝酸和蛋白质发生了颜色反应的缘故。

此外，蛋白质被灼烧时，产生具有烧焦羽毛的气味。

蛋白质是人和动物不可缺少的营养物质。
人们从食物摄取的蛋白质，在胃液中的胃蛋白酶和胰液中的胰蛋白酶的作用下，经过水解反应，生成氨基酸。
氨基酸被人体吸收后，再重新结合成人体所需要的各种蛋白质。

生命现象跟蛋白质有密切关系，没有蛋白质就没有生命。
1965 年我国科学家在世界上第一次用人工方法合成了具有生命活力的蛋白质一结晶牛胰岛素。
这对蛋白质和生命的研究做出了贡献。

蛋白质不仅具有生物学上的意义，在工业上也有广泛用途。
动物的毛和蚕丝的成分都是蛋白质，它们是重要的纺织原料。
动物的皮经过药剂鞣制后，其中所含蛋白质就变成不溶于水的，不易腐烂的物质，可以加工制成柔软坚韧的皮革。

动物胶的主要成分是蛋白质，是用骨和皮等熬煮而得的。
无色透明的动物胶叫做白明胶，是制造照相感光片和感光纸的原料。

牛奶里的蛋白质——酪素除做食品外，还能跟甲醛合成酪素塑料。

\subsection{酶}
酶是一类蛋白质，它们具有蛋白质的特性，如温度稍高或遇强酸、强碱、重金属盐类时能够变性等。
酶也有它们自己的特性，酶是生物制造出来的催化剂，催化许多有机化学反应和生物体中复杂的反应，例如氧化—还原、水解等。
酶脱离生物机体后仍具有活性。
前面所讲到的淀粉和纤维素在微生物作用下发酵变为葡萄糖，葡萄糖在微生物作用下发酵变为酒精，都有微生物产生的各种酶的作用。

在人和动植物的生理活动中，酶起着重要的作用，如食物中的淀粉为人们的唾液和胰液中含有的淀粉酶所水解。
食物中的脂肪为胰液和小肠液中的脂肪酶所水解，食物中的蛋白质为胃液和胰液中的蛋白酶所水解。
\begin{Experiment}
在一个试管里注入 \qty{10}{ml} 新配制的 1\% 淀粉溶液，加入约 \qty{1}{ml} 的唾液。把它们混和均匀，让混和液继续保持在室温条件下。取一个试管，放入混和液 5 滴，用碘液试验。同时用未加唾液的淀粉溶液作对比试验。
\end{Experiment}

分析上述实验时，要跟淀粉用强酸催化发生水解的实验作对比。
可以看出，用淀粉酶水解淀粉比用强酸水解，一是条件温和，即不需加热；二是反应快，效率高。
酶催化还有另一个特点就是它的专一性，如淀粉酶只催化淀粉的水解。

人们现在已经知道的酶有一千种以上。
工业上大量使用的酶多数是通过微生物发酵制得的，并且已经有许多种酶制成了晶体。
酶已得到广泛的应用，如淀粉酶应用于食品、发酵、纺织、制药等工业；蛋白酶用于医药、制革工业等；脂肪酶用来使脂肪水解、羊毛脱脂等。
酶也用于化学工业制造多种有机溶剂和试剂，如柠檬酸、丙酮、丁醇等。

\begin{Practice}[习题]
\begin{question}
  \item 什么是氨基酸？它有什么特性？
  \item 某种蛋白质含有 0.64\% 的硫，它的分子里只含有两个硫原子，计算这种蛋白质的式量。
  \item 在三个试管里分别盛有蛋白质、淀粉和肥皂的溶液，怎样鉴别它们？
  \item 根据蛋白质的性质，回答下列问题：
  \begin{tasks}[label=(\arabic*)]
    \task 为什么用煮沸的方法可以使医疗器械消毒？
    \task 为什么硫酸铜或氯化汞溶液能杀菌？
    \task 为什么铜、汞、铅等重金属盐类对人畜有毒？
    \task 误服重金属盐，为什么可以服用大量牛乳、蛋清或豆浆解毒？
    \task 用灼烧的方法为什么可以区别毛织物或棉织物？
  \end{tasks}
\end{question}
\end{Practice}

\section*{内容提要}
\addcontentsline{toc}{section}{内容提要}\setcounter{subsection}{0}
\subsection*{一、糖类}
\begin{enumerate}[1.]
  \item 单糖是多羟基醛或多羟基酮，如葡萄糖、果糖等。单糖不能水解成更简单的糖. 葡萄糖具有还原性，它是重要的营养物质。
  \item 低聚糖水解后能生成两个或几个分子的单糖，低聚糖中重要的是二糖，如蔗糖、麦芽糖等。蔗糖不具有还原性，水解产物有还原性，它是重要的甜味食物。
  \item 多糖是由多个单糖分子按照一定方式通过脱去水分子结合而成的高分子化合物，如淀粉、纤维素等。淀粉、纤维素都不具有还原性。多糖能够水解成为单糖，水解产物有还原性。淀粉、纤维素在生活上、生产上都有重要用途。
\end{enumerate}
\subsection*{二、蛋白质}
\begin{enumerate}[1.]
  \item 氨基酸\quad 蛋白质经水解，生成 $\alpha$-氨基酸。氨基酸既显酸性，又显碱性。
  \item 蛋白质是由 $\alpha$-氨基酸通过肽键构成的高分子化合物。少量的盐可促进蛋白质溶解，多量的盐可使蛋白质从溶液中析出。蛋白质既显酸性，又显碱性。蛋白质受热或受重金属盐类作用即变性。蛋白质是生命的物质基础。
  \item 酶的成分是蛋白质，能催化有机体中或有机物的许多反应，它是有机催化剂。酶的催化作用具有专一性、效率高、反应条件温和等特性。酶的催化反应已经广泛地应用于生产上。
\end{enumerate}

\begin{Review}
\begin{question}
  \item 下列说法有没有错误？说明理由。
  \begin{tasks}[label=(\arabic*)]
    \task 脂肪、碳水化合物、蛋白质都是人类食物的主要成分，它们都是天然有机高分子化合物。
    \task 碳水化合物是碳的水化物。
    \task 鸡蛋白的溶液是一种溶液，但是它具有溶胶的性质。
    \task 二糖是含有两个碳原子的糖。
  \end{tasks}
  \item 蔗糖和麦芽糖是不是同分异构体？淀粉和纤维素是不是同分异构体？说明理由。
  \item 在由淀粉制葡萄糖的过程中，怎样判断水解反应处于
  \begin{enumerate*}[(1)]
    \item 淀粉尚未水解；
    \item 淀粉正在水解；
    \item 淀粉的水解已完成。
  \end{enumerate*}
  \item 火棉、醋酸纤维、人造丝、滤纸的成分有什么不同？
\end{question}
\end{Review}